\scanpage{207}\mspage{152}
$Y$ de $\underline{\mathrm{Liss}}\,S$. En fait, on n'utilise que la lissité de $X,X'$, pas celle de $Y$, et c'est immédiat, car \deleted{d'après le} si $X\longrightarrow X'$ est dans $W_S^d$, $X\times_S Y\longrightarrow X'\times_S Y^{\deleted{\unsure{S}}}$ est dans $W_Y^d$, \deleted{\unsure{cf. th. 1}\ \illegible{2}} donc aussi dans $W_S^d$ (cf th. 1 appliqué à $Y\longrightarrow S$). Cette remarque \deleted{\unsure{prouve}} \inserted{suggère} en même temps :

\uline{Corollaire 1} \inserted{(\uuline{Faux}\unsure{!})} La formule bifonctorielle $(\unsure{+|*})$ est valable pour $X$ \deleted{lisse,} dans $\underline{\mathrm{Liss}}\,S$, $Y$ quelconque dans $\mathrm{Cat}_{/S}$, et aussi dans le cas inverse.
% slanted
\margin{Cet énoncé est sûrement \uuline{faux} ! \unsure{Contrexemple|Contreexemple} p. 153. Cf cor. 1 rectifié page \unsure{154|155}\deleted{\unsure{155}}}

Dans le premier cas, considérons
\[
 Y\xrightarrow{i_Y}\Phi(Y),\quad\text{dans }W_S^d\text{ donc}
\]
\deleted{\unsure{aussi}} iso dans $(\mathrm{Cat}_{/S})(W_S^d)^{-1}$, on aura donc
\[
\begin{aligned}
 \mathrm{hot}_{S,\unsure{W|w}}(X)\times_S\mathrm{hot}_{S,W}(Y)&\xrightarrow{\sim}\mathrm{hot}_{S,W}(X)\times_S\mathrm{hot}_{S,W}(\Phi(Y))\\
 &\overset{\substack{\text{car }X\ \text{\uline{et}}\ \Phi(Y)\ \text{lisses sur }S}}{\xrightarrow{\sim}}\mathrm{hot}_{S,W}(X\times_S\Phi(Y)),
\end{aligned}
\]
\deleted{\unsure{or}} je dis que
\[
 X\times_S Y\xrightarrow{X\times_S i_Y}X\times_S\Phi(Y)
\]
est encore dans $W_S^d$. Il suffit de voir que c'est dans $W_X^d$ (cf th. 1), \deleted{\unsure{donc}} \inserted{et pour ceci,} que le chgt. de base lisse $X\longrightarrow S$ induit $\mathrm{Cat}_{/S}\longrightarrow\mathrm{Cat}_{/\unsure{x|X}}$ compatible avec les localiseurs $W_S^d,W_X^d$. \deleted{Ce point étant admis, \unsure{donc}\ \illegible{1}} \inserted{Mais c'est \uline{faux} (prendre $X=\{s\}$) :}
% archive sheet number: 206
